\section{Conclusion}
\label{sec:conclusion}

We present the first empirical study of single-annotator RLHF as a diagnostic tool for understanding verbosity and blandness in language models. By training DPO models on individual simulated annotators with diverse length preferences, we show that the verbosity problem in standard RLHF is largely a preference-aggregation artifact: aggregate training amplifies verbosity by 21\% over the SFT base model, while a brevity-preferring annotator reduces it by 30\% relative to aggregate ($p{=}0.042$). However, the ``all AI sounds the same'' problem persists across all conditions---an LLM judge rates distinctiveness at only 1.3--1.6 out of 5 regardless of whether the model uses aggregate, individual, or no RLHF at all. We conclude that blandness is rooted in pretraining and SFT, not in preference optimization.

For practitioners, our results suggest that personalized RLHF is a viable tool for controlling verbosity but not for creating distinctive AI voices. For researchers, the decomposition points toward pretraining and SFT as the stages most in need of intervention if the goal is genuinely diverse model outputs. Future work should extend this analysis to larger models, real human annotators, and other dimensions of the blandness problem beyond verbosity, including hedging language, formulaic structure, and excessive use of lists.
